\chapter{Teoretická východiska}
\label{chap:teorie}

Tato kapitola analyzuje klíčové technologie, architektonické principy a bezpečnostní standardy nezbytné pro návrh moderního CLI nástroje v prostředí OS Linux. Text postupuje od obecné problematiky příkazové řádky přes síťové protokoly až po specifika implementačních platforem .NET a Go.

\section{Rozhraní příkazové řádky}
\label{sec:cli_teorie}
Rozhraní příkazové řádky (CLI – Command Line Interface) představuje jeden z nejstarších, avšak stále nejvýkonnějších způsobů interakce uživatele s počítačem. Navzdory masivnímu rozšíření grafických uživatelských rozhraní (GUI) v 90. letech 20. století zůstává CLI nepostradatelným nástrojem pro systémové administrátory, vývojáře a pokročilé uživatele, zejména v prostředí operačních systémů typu Unix a Linux.

\subsection{Role CLI v moderní správě systémů}
Hlavním důvodem preference příkazové řádky v profesionálním prostředí je efektivita, rychlost a snadná automatizace. Zatímco grafické rozhraní je navrženo pro intuitivní objevování funkcí pomocí vizuálních prvků, CLI se zaměřuje na precizní a opakovatelné provádění úkolů. Administrátor může pomocí jediného příkazu nebo krátkého skriptu provést operaci na stovkách serverů současně, což je v rámci GUI prakticky neproveditelné.

Dalším kritickým faktorem je vzdálená správa. Protokol SSH (Secure Shell), který je de facto standardem pro správu linuxových serverů, je primárně textově orientovaný. Přenos textových příkazů a jejich výstupů vyžaduje minimální síťovou propustnost, což umožňuje spolehlivou správu systémů i přes nestabilní nebo pomalé připojení. V prostředí cloudu a kontejnerizace (např. Docker, Kubernetes), kde systémy často nemají žádné grafické výstupní zařízení (headless režim), je CLI jedinou možností interakce.

\subsection{Principy návrhu a standardy CLI aplikací}
Kvalitní CLI aplikace se řídí souborem ustálených principů, které zajišťují jejich předvídatelnost a vzájemnou spolupráci. Eric S. Raymond ve své knize \textit{The Art of Unix Programming} definuje tzv. „Unixovou filosofii“, jejímž základem je tvorba malých, jednoúčelových nástrojů, které spolu efektivně komunikují \cite{raymond_unix}. Tento přístup umožňuje skládat komplexní operace z jednoduchých komponent pomocí mechanismu rour (pipes).

Základním stavebním kamenem interakce v CLI jsou standardní datové proudy:
\begin{itemize}
    \item \textbf{Standardní vstup (stdin):} Zdroj dat pro aplikaci, typicky klávesnice nebo výstup jiného programu.
    \item \textbf{Standardní výstup (stdout):} Primární kanál pro výstup dat, typicky terminál.
    \item \textbf{Standardní chybový výstup (stderr):} Oddělený kanál pro chybová hlášení a diagnostiku, což umožňuje uživateli přesměrovat data do souboru, zatímco chyby stále vidí na obrazovce.
\end{itemize}

Důležitou roli hraje také syntaxe příkazů. Moderní nástroje se snaží dodržovat standardy jako POSIX nebo GNU konvence pro argumenty a přepínače (flags). Krátké přepínače (např. \texttt{-v}) jsou určeny pro rychlé psaní, zatímco dlouhé varianty (např. \texttt{--verbose}) zvyšují čitelnost skriptů a dokumentace.

\subsection{CLI versus TUI}
V rámci terminálu se můžeme setkat také s textovým uživatelským rozhraním (TUI – Text User Interface). Na rozdíl od klasického CLI, které vypisuje řádky textu v sekvenčním pořadí, TUI využívá celou plochu terminálu k vykreslování vizuálních prvků, jako jsou okna, menu, tabulky nebo interaktivní grafy, často s využitím knihoven jako \texttt{ncurses} nebo moderních frameworků v jazyce Go (např. \texttt{Bubble Tea}).

Hlavní rozdíl spočívá v komplexitě a způsobu interakce. CLI je ideální pro jednorázové příkazy a automatizaci v nekonečném proudu (batch processing). TUI je naproti tomu vhodnější pro interaktivní nástroje, kde uživatel potřebuje neustálý přehled o stavu systému nebo kde je navigace v datech pomocí klávesových zkratek efektivnější než vypisování parametrů (např. textové editory, správci souborů nebo dashboardy pro monitoring).

\section{Architektura VPN a principy Zero Trust}
\label{sec:vpn_zero_trust}
Tradiční pojetí zabezpečení sítě, založené na ochraně perimetru, prochází v posledních letech zásadní transformací. Tato sekce popisuje evoluci od klasických VPN tunelů k moderní architektuře Zero Trust Network Access (ZTNA).

\subsection{Tradiční VPN a perimetrová bezpečnost}
Virtuální privátní síť (VPN) je technologie, která umožňuje vytvořit bezpečné, šifrované spojení (tunel) mezi dvěma body v rámci nezabezpečené sítě, jako je internet. Princip tunelování spočívá v zapouzdření (enkapsulaci) síťových paketů do jiného protokolu, který zajišťuje integritu a důvěrnost přenášených dat.

Historicky se organizace spoléhaly na model „hrad a příkop“ (Castle-and-Moat). V tomto modelu je vnitřní síť považována za důvěryhodnou zónu, která je od vnějšího světa oddělena firewallem (příkopem). VPN v tomto kontextu slouží jako „padací most“, který umožňuje vzdáleným uživatelům vstoupit do vnitřní sítě. Jakmile však uživatel získá přístup, má často široký dosah k mnoha zdrojům v rámci sítě. Tento přístup trpí kritickou slabinou: pokud útočník kompromituje jediné zařízení uvnitř perimetru, může se v síti volně pohybovat (laterální pohyb) a napadat další systémy, protože vnitřní síť mu implicitně důvěřuje.

\subsection{Zero Trust Network Access (ZTNA) a GoodAccess}
Filozofie Zero Trust (nulová důvěra) tento model opouští a vychází z principu „nikdy nedůvěřuj, vždy prověřuj“ (Never trust, always verify). V architektuře ZTNA není důvěra nikdy udělena na základě fyzického umístění uživatele nebo jeho IP adresy. Každý pokus o přístup k jakémukoliv zdroji musí být individuálně autentizován, autorizován a šifrován.

Klíčovým prvkem ZTNA je Software Defined Perimeter (SDP). SDP vytváří dynamickou, virtuální hranici kolem aplikací a služeb, čímž je činí „neviditelnými“ pro veřejný internet i pro neautorizované uživatele v lokální síti. GoodAccess implementuje tyto principy formou cloudové platformy, která nahrazuje tradiční VPN brány \cite{goodaccess_doc}. Místo statických IP adres využívá identitu uživatele (často propojenou s SSO/OIDC) a kontext zařízení (tzv. Device Posture) k dynamickému přidělování přístupových práv k jednotlivým systémům \cite{goodaccess_architecture}. Tím se radikálně snižuje plocha možného útoku, protože uživatel vidí pouze ty zdroje, ke kterým má explicitní oprávnění.

\subsection{Srovnání protokolů WireGuard a OpenVPN}
Pro realizaci samotného šifrovaného spojení na platformě Linux využívá GoodAccess dva dominantní protokoly: OpenVPN a WireGuard. Každý z nich představuje jinou generaci technologií a liší se především architekturou a flexibilitou správy.

OpenVPN je dlouholetým standardem v oblasti VPN. Je postaven na knihovně OpenSSL a nabízí extrémně široké možnosti konfigurace, včetně podpory různých šifrovacích algoritmů a transportních protokolů (UDP i TCP). Tato univerzalita je však vykoupena vysokou komplexitou zdrojového kódu (stovky tisíc řádků), což ztěžuje bezpečnostní audity a vede k vyšší režii při zpracování paketů, zejména v uživatelském prostoru.

Naproti tomu WireGuard je moderní protokol navržený s důrazem na jednoduchost, rychlost a moderní kryptografii \cite{wireguard_paper}. S rozsahem přibližně 4 000 řádků kódu je WireGuard snadno auditovatelný a díky implementaci přímo v jádře Linuxu dosahuje výrazně vyšší propustnosti a nižší latence než OpenVPN. Z pohledu správy systému se WireGuard chová jako standardní síťové rozhraní, což zjednodušuje konfiguraci směrování a integraci do systémových nástrojů. Pro dynamické enterprise prostředí, které GoodAccess obsluhuje, je WireGuard ideální volbou díky rychlému navazování spojení (handshake) a vysoké odolnosti vůči změnám síťového prostředí.

\section{Architektura OS Linux a správa systému}
\label{sec:linux_architektura}
Pro pochopení interakce mezi CLI klientem a systémovými prostředky je nutné definovat hranice mezi uživatelským prostorem a jádrem.

\subsection{Jádro, Netlink a virtuální FS}
Operační systém Linux striktně odděluje privilegovaný prostor jádra (Kernel Space) od uživatelského prostoru (User Space). Pro aplikace typu VPN klienta je klíčová schopnost konfigurovat síťové parametry, které spravuje právě jádro. K tomuto účelu Linux poskytuje několik rozhraní.

Virtuální souborové systémy \texttt{/proc} a \texttt{/sys} umožňují aplikacím číst stav systému a měnit některé parametry pouhým zápisem do souboru (např. povolení IP forwardingu). Nicméně pro komplexnější operace, jako je vytváření síťových rozhraní nebo správa směrovacích tabulek, se využívá rozhraní \texttt{Netlink}. Jak uvádí Ward \cite{ward_linux}, \texttt{Netlink} je socketové rozhraní sloužící ke komunikaci mezi jádrem a uživatelskými procesy. V této práci je \texttt{Netlink} využíván nepřímo skrze systémové nástroje nebo knihovny v .NET, které zajišťují nízkoúrovňovou konfiguraci WireGuard rozhraní.

\subsection{Správa služeb a Systemd}
Moderní linuxové distribuce využívají jako inicializační systém a správce služeb \texttt{systemd}. Pro aplikaci GoodAccess CLI je \texttt{systemd} zásadní, protože zajižduje běh backendové služby (\texttt{CLIService}) jako systémového démona, který je spouštěn při startu systému a automaticky restartován v případě chyby. \texttt{systemd} také umožňuje definovat závislosti služeb, což zaručuje, že se VPN služba pokusí o navázání spojení až ve chvíli, kdy je dostupné síťové rozhraní. V rámci této práce jsou vytvořeny \textit{unit} soubory, které definují životní cyklus služby a její integraci do systému.

\subsection{Bezpečné ukládání tajemství (Secure Storage)}
Kritickým aspektem VPN klienta je ochrana privátních klíčů a tokenů. Linux standardně využívá \textit{Secret Service API} komunikující přes DBus, což umožňuje aplikacím ukládat tajemství do klíčenky (Keyring) spravované systémem, nikoliv do prostých textových souborů.

\section{Implementační technologie}
\label{sec:technologie}
Architektura řešení je rozdělena na privilegovanou službu (Daemon) a uživatelské rozhraní (Client).

\subsection{Backend: Platforma .NET}
Pro implementaci systémové služby (Daemon) byla zvolena platforma .NET 8. Tato volba vychází především z potřeby kontinuity s existující infrastrukturou GoodAccess, která je na technologiích Microsoft .NET postavena. .NET na Linuxu dnes nabízí vysoký výkon a plnou podporu pro systémové programování, včetně nativní interakce s nízkoúrovňovými API systému přes P/Invoke nebo knihovny typu Microsoft.Extensions.Hosting pro běh služeb. Modul \texttt{CLIService} tak může sdílet klíčovou logiku pro správu tokenů, certifikátů a stavu tunelu se stávajícím grafickým klientem, což zaručuje konzistentní chování aplikace napříč různými rozhraními.

\subsection{Frontend: Jazyk Go a TUI}
Pro klientskou část (CLI frontend) byl zvolen jazyk Go. Go se v posledních letech stal de facto standardem pro vývoj cloud-native a systémových nástrojů v ekosystému Linux (např. Docker, Kubernetes, Terraform). Hlavní výhodou je schopnost kompilace do jediné staticky linkované binárky bez nutnosti instalace runtime prostředí, což je pro nástroj příkazové řádky ideální. Go nabízí vynikající knihovny pro tvorbu moderních terminálových rozhraní – v této práci je využita knihovna \texttt{Cobra} pro definici struktury příkazů a framework \texttt{Bubble Tea} pro reaktivní zobrazení stavu a interaktivní prvky.

\subsection{Meziprocesová komunikace (IPC)}
V architektuře, kde je logika aplikace oddělena od uživatelského rozhraní do samostatných procesů, hraje klíčovou roli meziprocesová komunikace (IPC – Inter-Process Communication). Pro potřeby této práce byla zvolena technologie Unix Domain Sockets (UDS).

Unix Domain Sockets představují datový komunikační koncový bod pro výměnu dat mezi procesy běžícími na stejném hostitelském operačním systému. Na rozdíl od síťových socketů (TCP/IP), které jsou určeny pro komunikaci přes síť, UDS nevyužívají síťový zásobník (network stack), což výrazně snižuje režii a zvyšuje propustnost komunikace.

Z pohledu bezpečnosti, která je u VPN klienta kritická, nabízejí UDS zásadní výhodu v podobě využití standardních souborových oprávnění systému Linux. Přístup k socketu lze omezit na konkrétního uživatele nebo skupinu, čímž se zamezí neoprávněným aplikacím v komunikaci se systémovou službou. V navržené architektuře slouží UDS jako transportní vrstva pro protokol gRPC, který zajišťuje silně typované rozhraní pro příkazy jako autentizace, připojení nebo monitoring stavu.

\section{Distribuce softwaru}
\label{sec:distribuce}
Zadání práce vyžaduje přípravu balíčků pro distribuce Debian a Fedora.
% Zde popiš .deb a .rpm formáty

\section{Metodika vývoje}
\label{sec:metodika}
% Zde použij svůj text o XP a Beckovi
